% =======================
% Section 4: Coding Agents
% =======================
\section{Coding Agents}

% ---- Slide 12: Coding Agents ----
\begin{frame}{Coding Agents and Long-Horizon Autonomous Agents}
\begin{itemize}
  \item \textbf{Why code is a good domain for agents:}
  \begin{itemize}
    \item The environment is structured and well-defined
    \item Tools are explicit: shell, file system, test runner
    \item Success is verifiable — tests pass or fail
    \item Iteration is natural — write, run, observe, fix
  \end{itemize}
  \vspace{0.3em}
  \item Current examples: \textbf{Codex CLI} (OpenAI), \textbf{Claude Code} (Anthropic)
  \item For longer-horizon autonomy: \textbf{OpenClaw} illustrates the direction toward persistent autonomous agents
\end{itemize}
\vspace{0.4em}
\textbf{A spectrum of autonomy:}
\vspace{0.3em}
\begin{center}
\begin{tikzpicture}[
  box/.style={draw, rounded corners=3pt, fill=blue!10,
              minimum width=2.5cm, minimum height=0.6cm,
              align=center, font=\scriptsize},
  every node/.style={font=\scriptsize}
]
  \node[box]                         (a1) at (0,0)   {Assistant};
  \node[box, fill=blue!18]           (a2) at (2.7,0) {Tool-using\\[-2pt]Agent};
  \node[box, fill=orange!18]         (a3) at (5.4,0) {Coding\\[-2pt]Agent};
  \node[box, fill=red!12, text width=2.2cm] (a4) at (8.1,0) {Long-horizon\\[-2pt]Autonomous};
  \draw[-{Stealth}] (a1) -- (a2);
  \draw[-{Stealth}] (a2) -- (a3);
  \draw[-{Stealth}] (a3) -- (a4);
\end{tikzpicture}
\end{center}
% Speaker notes:
% Verifiability is key: agents work well when "did it work?" has a clear answer.
% That's why code, math, and data pipelines are natural agent domains.
% Reliability, security, and oversight remain major challenges at the long-horizon end.
\end{frame}

% ---- Slide 13: Codex — The Harness and Prompt Assembly ----
\begin{frame}{Codex: The Harness and Prompt Assembly}
\begin{columns}[T]
\begin{column}{0.52\textwidth}
\textbf{Key insight}: Codex separates the \emph{model} from the \emph{harness}
\begin{itemize}
  \item The \textbf{harness} is the orchestration layer that assembles prompts, drives the loop, executes tools, and manages state
  \item ``LLM + tools'' elides exactly this layer — the harness is the agent
\end{itemize}
\vspace{0.3em}
\textbf{Prompt assembly is dynamic} — the prompt the user thinks they send is not what the model sees. The harness assembles it in layers:
\begin{enumerate}
  \item System instructions \textit{(model-specific, bundled in CLI)}
  \item Sandbox / permissions context
  \item Developer instructions
  \item User instructions \textit{(from \texttt{AGENTS.md} files / skills)}
  \item Environment context \textit{(cwd, shell)}
  \item User message
\end{enumerate}
\vspace{0.2em}
\small Native tools, Responses API tools, and MCP tools all appear \emph{identical} to the model in the tools array.
\end{column}
\begin{column}{0.44\textwidth}
\begin{tikzpicture}[
  box/.style={draw, rounded corners=3pt, fill=blue!10,
              minimum width=3.0cm, minimum height=0.55cm,
              align=center, font=\scriptsize},
  every node/.style={font=\scriptsize}
]
  \node[box, fill=gray!12]   (sys)  at (0, 4.2) {1. System instructions};
  \node[box, fill=red!10]    (sand) at (0, 3.55) {2. Sandbox / permissions};
  \node[box, fill=orange!12] (dev)  at (0, 2.9) {3. Developer instructions};
  \node[box, fill=yellow!15] (agmd) at (0, 2.25) {4. AGENTS.md / skills};
  \node[box, fill=blue!12]   (env)  at (0, 1.6) {5. Environment context};
  \node[box, fill=green!15]  (usr)  at (0, 0.95) {6. User message};
  \node[box, fill=purple!15] (inf)  at (0, 0.2) {Model inference};

  \draw[-{Stealth}] (sys)  -- (sand);
  \draw[-{Stealth}] (sand) -- (dev);
  \draw[-{Stealth}] (dev)  -- (agmd);
  \draw[-{Stealth}] (agmd) -- (env);
  \draw[-{Stealth}] (env)  -- (usr);
  \draw[-{Stealth}] (usr)  -- (inf);

  \draw[dashed, gray, rounded corners=3pt]
    (-1.62, 4.5) rectangle (1.62, 0.6);
  \node[font=\scriptsize\bfseries, text=gray] at (0, 4.72) {Harness assembles};
\end{tikzpicture}
\end{column}
\end{columns}
% Speaker notes:
% Connect to Slide 2 (Anatomy): the harness IS the control + action + state layers made concrete.
% Skills appear here as AGENTS.md / skill files — this is the Codex instantiation of the concept.
% MCP safety reminder: MCP tools are NOT sandboxed by Codex.
\end{frame}

% ---- Slide 14: Codex — Managing the Growing Context ----
\begin{frame}{Codex: Managing the Growing Context}
\begin{columns}[T]
\begin{column}{0.52\textwidth}
\textbf{The prompt grows monotonically:}
\begin{itemize}
  \item Each turn appends to prior input — full conversation history always included
  \item A single turn can contain many iterations of inference + tool calls
  \item A turn ends only when the model emits an \emph{assistant message} (termination signal)
  \item This is \textbf{quadratic} in data sent — a real engineering problem
\end{itemize}
\vspace{0.3em}
\textbf{Two mitigations:}
\begin{itemize}
  \item \textbf{Prompt caching} — reuse computation from prior inference if the prefix is stable. Static content (instructions, tools) must come first; variable content last. \textit{Cache miss}: MCP tools enumerated in inconsistent order $\to$ expensive miss every turn.
  \item \textbf{Automatic compaction} — when token threshold exceeded, conversation is summarized and replaced with a compressed representation, preserving the model's latent understanding via encrypted content
\end{itemize}
\end{column}
\begin{column}{0.44\textwidth}
\begin{tikzpicture}[
  box/.style={draw, rounded corners=3pt, fill=blue!10,
              minimum width=2.8cm, minimum height=0.55cm,
              align=center, font=\scriptsize},
  every node/.style={font=\scriptsize}
]
  \node[box, fill=gray!12]   (t1) at (0, 4.2) {Turn 1: prompt};
  \node[box, fill=blue!10]   (i1) at (0, 3.5) {Inference + tool calls};
  \node[box, fill=green!12]  (a1) at (0, 2.8) {Assistant message $\to$ end turn};

  \node[box, fill=gray!20]   (t2) at (0, 2.0) {Turn 2: prompt\\(Turn 1 + new msg)};
  \node[box, fill=blue!10]   (i2) at (0, 1.2) {Inference + tool calls};
  \node[box, fill=green!12]  (a2) at (0, 0.5) {Assistant message $\to$ end turn};

  \node[box, fill=orange!20] (cp) at (0,-0.3) {Compaction\\(if threshold exceeded)};

  \draw[-{Stealth}] (t1) -- (i1);
  \draw[-{Stealth}] (i1) -- (a1);
  \draw[-{Stealth}] (a1) -- (t2);
  \draw[-{Stealth}] (t2) -- (i2);
  \draw[-{Stealth}] (i2) -- (a2);
  \draw[-{Stealth}, dashed] (a2) -- (cp);

  % Growing arrow
  \draw[{Stealth}-, thick, blue!40] (-1.7, 4.4) -- (-1.7, -0.0)
    node[midway, left, font=\scriptsize, blue!60] {\shortstack{growing\\prompt}};
\end{tikzpicture}
\end{column}
\end{columns}
% Speaker notes:
% Connect to Appendix B failure mode "excessive cost / latency" — now students see the mechanism.
% The MCP tool reordering bug is a perfect real example of how implementation details
% have direct cost impact in production.
% Key lesson: the agent loop is not just about reasoning — it is about context budget management.
\end{frame}
